\documentclass[../../main.tex]{subfiles}
\begin{document}
% SOURCE p=242; printed=227; object=ch07-radon-product-definition; region=middle
现在已清楚如何由 $X,Y$ 上的 Radon 测度 $\mu,\nu$ 得到 $X\times Y$ 上的 Radon 测度：由\cref{proposition:ch07-cc-product-integrals}，$I(f)=\int f\,\mathrm d(\mu\times\nu)$ 在 $C_c(X\times Y)$ 上定义了正线性泛函，所以 Riesz 表示定理确定一个 Radon 测度。
\begin{definition}[Radon 乘积]\label{definition:ch07-radon-product}
把上述由连续紧支撑函数的积分泛函确定的 Radon 测度称为 $\mu,\nu$ 的\textbf{Radon 乘积}，记作 $\mu\mathbin{\widehat\times}\nu$。即对每个 $f\in C_c(X\times Y)$，
\begin{align*}
\int f\,\mathrm d(\mu\mathbin{\widehat\times}\nu)=\int f\,\mathrm d(\mu\times\nu).
\end{align*}
\end{definition}
% SOURCE p=242; printed=227; object=ch07-radon-product-counterexample; region=middle
自然的问题是：$\mu\mathbin{\widehat\times}\nu$ 在 $\mathscr B_X\otimes\mathscr B_Y$ 上是否与 $\mu\times\nu$ 一致？一般而言，答案是否定的。可取 $X=\mathbb R$、$Y=\mathbb R_d$（赋予离散拓扑的实数集），$\mu$ 为 Lebesgue 测度、$\nu$ 为计数测度作为反例。不难看出此时 $\mathscr B_{X\times Y}=\mathscr B_X\otimes\mathscr B_Y$，而\cref{exercise:ch07-uncountable-lines,exercise:ch07-semifinite-radon}说明 $\mu\mathbin{\widehat\times}\nu$ 不半有限。原书接着断言：“$\mu\times\nu$ 是 $\mu\mathbin{\widehat\times}\nu$ 的半有限部分。”此句按本书普通乘积的定义需要订正，详见下面的注解。不过，仍有若干可用的结论，并且在 $\sigma$ 有限情形中，各种关系都能良好成立。以下沿用\cref{section:ch02-product-measures}中的 $x$ 截面和 $y$ 截面记号。
% SUPPLEMENT: id=v14-ch07-radon-product-definition-details; parent=ch07-radon-product-definition; kind=proof-detail+motivation+hypotheses; origin=independent-derivation
\begin{remark}
\textbf{构造的动机与唯一性范围（重构新增）。}普通乘积首先服务于可测矩形，Radon 乘积则首先服务于连续紧支撑测试。前命题确保积分有限、线性且保持非负，所以正 Riesz 定理给出存在性与 Radon 类中的唯一性。乘积空间 LCH 是应用该定理的拓扑条件；两个因子的 Radon 性使紧支撑测试可积。定义没有预先宣称两个测度在全部公共可测集上相等，只有上述测试积分相同；外正则化可能增加某些集合的质量。
\end{remark}
% END SUPPLEMENT: id=v14-ch07-radon-product-definition-details
% SUPPLEMENT: id=v14-ch07-radon-product-counterexample-details; parent=ch07-radon-product-counterexample; kind=omitted-proof+correction+motivation+hypotheses; origin=independent-derivation
\begin{remark}
\textbf{反例计算与原文半有限断言的订正（重构新增）。}在上述 $\mathbb R\times\mathbb R_d$ 中，取 $\mathbb R$ 的可数开基 $\{B_n\}$。任意开集 $U$ 可写成 $\bigcup_n(B_n\times C_n)$，其中 $C_n=\{y:B_n\subseteq U^y\}$ 是离散空间中的开集：每个 $(x,y)\in U$ 的横截面 $U^y$ 开，因而有 $B_n$ 包含 $x$ 且包含于该截面。于是每个开集都在乘积 $\sigma$ 代数中，两种 Borel 结构确实相同。

每个紧支撑连续函数只有有限个非零横截面，其普通乘积积分为 $\sum_y\int f(x,y)\,\mathrm dx$，故 Riesz 唯一性表明 Radon 乘积正是\cref{exercise:ch07-uncountable-lines}中的测度 $\rho$。该题已证明 $E=\{0\}\times\mathbb R_d$ 满足 $\rho(E)=\infty$，而每个紧子集质量为零，且 $\rho$ 的半有限部分由\cref{exercise:ch07-semifinite-radon}给为
\begin{align*}
\rho_0(H)=\sum_{y\in\mathbb R}m(H^y).
\end{align*}
所以 $\rho_0(E)=0$，确有 $\rho$ 不半有限。另一方面，普通乘积对矩形给 $(m\times\nu)(E)=m(\{0\})\nu(\mathbb R)=0$，其中使用 $0\cdot\infty=0$ 的定义约定，故它与 Radon 乘积不相等。

但普通乘积\emph{并不等于} $\rho_0$。为核验，取闭对角线 $D=\{(t,t):t\in\mathbb R\}$。由上段的 Borel 结构相等，$D$ 在乘积 $\sigma$ 代数内，且 $\rho_0(D)=\sum_y m(\{y\})=0$。反观\cref{definition:ch02-product-measure}中由矩形预测度诱导外测度的定义，$D$ 的普通乘积测度为无穷大。事实上，若一个可数矩形覆盖 $D\subseteq\bigcup_j(A_j\times B_j)$ 有有限总成本 $\sum_jm(A_j)\nu(B_j)$，则每个无限 $B_j$ 都必须满足 $m(A_j)=0$。把这些零测 $A_j$ 的可数并记为 $N$，把所有有限 $B_j$ 的并记为 $C$；$N$ 为零测集，$C$ 可数。取 $t\notin N\cup C$，则 $(t,t)$ 不在任何覆盖矩形中，矛盾。因此每个矩形覆盖成本都无穷，$(m\times\nu)(D)=\infty$。若起初采用有限不交矩形并的代数覆盖，将每项拆成有限个矩形就得到同一论证。

因此，按本书前面的普通乘积定义，正确说法是：Radon 乘积 $\rho$ 的半有限部分是上述截面求和测度 $\rho_0$；它通常不同于普通乘积，后者与 $\rho$ 也可不同。这个计算也与\cref{exercise:ch02-fubini-non-sigma-finite}的对角线现象一致。这里保留原断言并明确说明失败之处，不把“两个测度在紧支撑连续函数上相同”误用成对所有 Borel 集都相同。
\end{remark}
% END SUPPLEMENT: id=v14-ch07-radon-product-counterexample-details
% SOURCE p=242; printed=227; object=ch07-borel-sections; number=7.23; region=middle-lower
\begin{lemma}\label{lemma:ch07-borel-sections}
\begin{enumerate}
\item\label{item:ch07-borel-set-sections} 若 $E\in\mathscr B_{X\times Y}$，则每个 $x\in X$ 都有 $E_x\in\mathscr B_Y$，每个 $y\in Y$ 都有 $E^y\in\mathscr B_X$。
\item\label{item:ch07-borel-function-sections} 若 $f:X\times Y\to\mathbb C$ 为 Borel 可测函数，则每个 $f_x$ 是 $\mathscr B_Y$ 可测的，每个 $f^y$ 是 $\mathscr B_X$ 可测的。
\end{enumerate}
\end{lemma}
% SOURCE p=242; printed=227; object=ch07-borel-sections-proof; region=lower
\begin{proof}
所有截面 $E_x,E^y$ 分别对每个 $x,y$ 都 Borel 的集合 $E\subseteq X\times Y$ 构成一个 $\sigma$ 代数。它包含全部开集：若 $E$ 开，则截面是它在连续映射 $y'\mapsto(x,y')$、$x'\mapsto(x',y)$ 下的原像，所以开。因此该 $\sigma$ 代数包含 $\mathscr B_{X\times Y}$，证明第一项。第二项来自
\begin{align*}
(f_x)^{-1}(A)=(f^{-1}(A))_x,\qquad
(f^y)^{-1}(A)=(f^{-1}(A))^y.
\end{align*}
\end{proof}
% SUPPLEMENT: id=v14-ch07-borel-sections-details; parent=ch07-borel-sections; kind=proof-detail+motivation+hypotheses; origin=independent-derivation
\begin{remark}
\textbf{全部截面的逐项核验（重构新增）。}全空间的截面为整个因子；$(E^c)_x=Y\setminus E_x$，$(\bigcup_nE_n)_x=\bigcup_n(E_n)_x$，所以符合截面 Borel 条件的集合对补和可数并封闭。横截面的两个公式同样成立，因此确为 $\sigma$ 代数。对函数取任意复平面的开集 $A$，函数 Borel 性与第一项使两个原像公式右端 Borel，给出可测性；非负扩展实值函数同样适用。

动机是：虽然一个 Borel 集未必在乘积 $\sigma$ 代数内，固定一个坐标却仍然是连续操作，因此每个截面都有可测性。本引理不要求测度或可数性，甚至不需要 LCH。它尚未断言截面\emph{测度}随另一变量可测，这个更强结论需要后面的论证。
\end{remark}
% END SUPPLEMENT: id=v14-ch07-borel-sections-details
% SOURCE p=242; printed=227; object=ch07-continuous-partial-integrals; number=7.24; region=bottom
\begin{lemma}\label{lemma:ch07-continuous-partial-integrals}
若 $f\in C_c(X\times Y)$，$\mu,\nu$ 分别为 $X,Y$ 上的 Radon 测度，则函数 $x\mapsto\int f_x\,\mathrm d\nu$ 和 $y\mapsto\int f^y\,\mathrm d\mu$ 连续。
\end{lemma}
% SOURCE p=242,243; printed=227-228; object=ch07-continuous-partial-integrals-proof; region=bottom-top
\begin{proof}
只写 $f_x$ 的证明。只需说明，对任意 $x_0\in X$、$\varepsilon>0$，存在 $x_0$ 的邻域 $U$，使 $x\in U$ 时 $\|f_x-f_{x_0}\|_u<\varepsilon$，因为这时
\begin{align*}
\left|\int(f_x-f_{x_0})\,\mathrm d\nu\right|
\leqslant\varepsilon\nu(\pi_Y(\operatorname{supp}f)).
\end{align*}
对每个 $y\in\pi_Y(\operatorname{supp}f)$，存在 $x_0,y$ 的邻域 $U_y,V_y$，使 $(x,z)\in U_y\times V_y$ 时 $|f(x_0,y)-f(x,z)|<\varepsilon/2$。取 $\pi_Y(\operatorname{supp}f)$ 的有限子覆盖 $V_{y_1},\ldots,V_{y_m}$，令 $U=\bigcap_{j=1}^mU_{y_j}$；剩余细节原书留给读者。
\end{proof}
% SUPPLEMENT: id=v14-ch07-continuous-partial-integrals-details; parent=ch07-continuous-partial-integrals; kind=omitted-proof+motivation+hypotheses; origin=independent-derivation
\begin{remark}
\textbf{有限覆盖后的细节与另一变量（重构新增）。}记 $K_Y=\pi_Y(\operatorname{supp}f)$。若 $z\in K_Y$，选 $j$ 使 $z\in V_{y_j}$。对 $x\in U$，$x,x_0$ 都属于 $U_{y_j}$，于是
\begin{align*}
|f(x,z)-f(x_0,z)|
\leqslant|f(x,z)-f(x_0,y_j)|+|f(x_0,y_j)-f(x_0,z)|<\varepsilon.
\end{align*}
若 $z\notin K_Y$，两个函数值都为零。因此所有 $z$ 上的差小于 $\varepsilon$，其上确界至多 $\varepsilon$，足以得到积分误差至多 $\varepsilon\nu(K_Y)$。严格的上确界 $<\varepsilon$ 可把构造中的 $\varepsilon/2$ 先改成 $\varepsilon/4$，得到上确界至多 $\varepsilon/2$。为使积分误差小于给定 $\delta>0$，取振幅预算 $\varepsilon=\delta/(1+\nu(K_Y))$ 即可；$\nu(K_Y)=0$ 时积分恒为零也包括在内。

交换 $X,Y$，对 $y_0$ 及紧集 $K_X=\pi_X(\operatorname{supp}f)$，以连续性取 $V_x\ni y_0,U_x\ni x$ 使 $|f(z,y)-f(x,y_0)|<\varepsilon/4$；$K_X$ 上取有限 $U_x$ 子覆盖并交相应 $V_x$。对每个 $z\in K_X$ 使用同一三角不等式，$K_X$ 外函数为零，得 $\|f^y-f^{y_0}\|_u\leqslant\varepsilon/2$。乘以有限的 $\mu(K_X)$，证明另一积分连续。

事实上，这两个连续函数分别在 $K_X,K_Y$ 外为零，故各自属于 $C_c$。动机是把单点连续性通过紧覆盖升级为对整个积分变量的一致控制；紧支撑提供有限覆盖及有限质量，Radon 的紧集有限性随后把振幅控制变成积分控制。
\end{remark}
% END SUPPLEMENT: id=v14-ch07-continuous-partial-integrals-details
\subfile{section-04-fubini.tex}
\end{document}
